\section{EXPERIMENTS}\label{experiments}

This section states the workload set, platforms, and
experiment-specific design choices. Shared timing, uncertainty,
and surface-validation rules are defined once in
Section~\ref{methods-benchmarking}. The three experiments form
a decision pipeline. RQ1 decides which GPU evaluation mode to
use and exports it as \texttt{rq3\_config}. RQ2 decides
whether heterogeneous scheduling is worth considering and on
which workload/platform regimes yield positive evidence. RQ3
decides whether calibration reuse amortizes only for regimes
with positive evidence in RQ2. Each stage depends on the
previous result, so a null or negative result upstream narrows
the choices available downstream
(Sections~\ref{experiment-1-algorithm}--\ref{experiment-3-cache}).

\subsection{Shared Experimental
Setup}\label{workload-set-and-platforms}

\textbf{Platforms.} Table~\ref{tab:platforms} gives the two
hardware settings: Apple~M4 with unified memory and ten CPU
cores, and NVIDIA~Tesla~T4 with discrete VRAM, higher GPU
throughput, and two CPU cores. The unified-vs-discrete
memory axis and the CPU-core count separate the RQ2 operating
points and bound the vmap chunk memory budget in RQ1.

\begin{table}[t]
  \caption{Platform Characterization}
  \label{tab:platforms}
  \centering
  \begin{tabular}{lll}
    \toprule
    Property & Apple M4 & NVIDIA Tesla T4 \\
    \midrule
    GPU & M4 integrated & Tesla T4 \\
    VRAM & 16\,GiB (unified) & 14.6\,GiB (discrete) \\
    CPU cores available & 10 & 2 \\
    PyTorch backend & \texttt{mps} & \texttt{cuda} \\
    \bottomrule
  \end{tabular}
\end{table}

\textbf{Workload set.} A workload is a (dataset, model family, loss function) triple. Four workloads cover
distinct computation patterns relevant to the RQ1 transform
applicability claim (Table~\ref{tab:workloads}):
\texttt{cifar10\_resnet20} for 2D convolution (the
workload for which CUDA convolution kernels are
designed~\cite{nvidia_cudnn_convolution_guide}),
\texttt{cifar10\_row\_gru}~\cite{cho2014gru,chung2014gru}
for recurrent row dependencies (no kernel fusion),
\texttt{mnist\_mlp} for MLP image classification, and
\texttt{california\_mlp} for
tabular regression. The MLP pair differs in input dimensionality and
per-sample compute cost, the axes that should impact vmap behavior. 
These span enough architectures
for workload-dependent outcomes to emerge. RQ1 and RQ2 use
all four; RQ3 narrows to one case selected by filtering
Experiment~2 for hybrid affinity
(Section~\ref{experiment-3-cache}).

\begin{table}[t]
  \caption{Workload set.}
  \label{tab:workloads}
  \centering
  \begin{tabular}{llll}
    \toprule
    Name & Modality & Architecture & Dataset \\
    \midrule
    \texttt{cifar10\_resnet20} & image & 2D conv (ResNet-20) & CIFAR-10 \\
    \texttt{cifar10\_row\_gru} & sequence & row-wise GRU & CIFAR-10 \\
    \texttt{mnist\_mlp} & dense image & MLP & MNIST \\
    \texttt{california\_mlp} & tabular & MLP & California Housing \\
    \bottomrule
  \end{tabular}
\end{table}

\textbf{Grid and data subset.} RQ1 and RQ2 use an
$8\times 8$ filter-normalized grid at unit perturbation
scale, 1024 evaluation samples per point, and batch size
$64$. Unit scale matches the upstream default for
filter-normalized
weights~\cite{li2018losslandscape,goldstein_loss_landscape_repo}.
The $8\times 8$ resolution yields 64 grid-point tasks per
run: enough to populate the $K\in\{32,64\}$ vmap chunks in
RQ1 and to fill the shared queue in RQ2, while keeping the
full sweep feasible on both platforms. Resolution sets the
task count, not the per-point evaluation semantics, so the
transform and scheduler verdicts do not depend on it. RQ3 uses a
$20\times 20$ session grid matching LossLens's default
resolution~\cite{xie2025losslens}, plus a calibration probe
grid of side $m$ with $m^2\ge 4(1+p_{\max})$
(Section~\ref{methods-calibration}).

\textbf{Repeats and trial order.} RQ1 and RQ2 use $R=5$
paired repeats ($df=R-1=4$, the minimum for the paired
t-CI). The paired design removes machine variance between runs,
so a small $R$ resolves the per-candidate verdicts across the
full $4\times 2$ workload/platform sweep at feasible total
cost; larger $R$ tightens the CI with diminishing
returns~\cite{kalibera2013rigorous}. Within each repeat the candidates run sequentially
under a cyclic rotation indexed by the repeat number, so no
candidate consistently runs after baseline and inherits its
warm cache
state~\cite{kalibera2013rigorous}. RQ3 takes a single
$T_{\text{session}}$ measurement per (workload, platform,
condition) and is descriptive
(Section~\ref{methods-benchmarking}).

\textbf{Conditions.} \texttt{vanilla} denotes the canonical
GPU-only path defined in
Section~\ref{methods-canonical-algorithm}. RQ1 pairs each
candidate from Section~\ref{methods-redesign} against
\texttt{vanilla}; RQ2 pairs \texttt{hybrid} against
\texttt{vanilla} under matched GPU slowdown; RQ3 pairs
\texttt{cached\_composed} against \texttt{vanilla} run
once per checkpoint.

\subsection{Algorithm-Level Optimization
(RQ1)}\label{experiment-1-algorithm}

\textbf{Question.} Do PyTorch's stateless evaluation
primitives reduce $T_{\text{grid}}$, and which workload
patterns gain or regress?

\textbf{Hypothesis.} Verdicts depend on architecture. We
predict supported improvement on the dense MLP rows (the
closest match to the \texttt{vmap} ensembling pattern), an
unresolved or smaller effect on \texttt{cifar10\_resnet20}
(its convolution kernels already saturate the GPU at this
batch size), and an unresolved or supported regression on
\texttt{cifar10\_row\_gru} (its inner row-by-row dependence
is not removable by batching perturbations).

\textbf{Design.} The candidates from
Section~\ref{methods-redesign} run on each workload/platform
pair under the shared setup. \texttt{vmap}-bearing
candidates sweep $K\in\{32,64\}$ as a memory-vs-throughput
sweep: the smaller chunk bounds memory and isolates the
dispatch-amortization gain, the larger chunk tests whether
higher per-chunk throughput exceeds the added memory
pressure. The composed candidate is compared against the
better individual transform via $q_{\text{compose}}$
(Section~\ref{methods-redesign}).

\subsection{Hybrid Scheduler Applicability
(RQ2)}\label{experiment-2-hybrid}

\textbf{Question.} Does heterogeneous CPU+GPU scheduling
reduce $T_{\text{grid}}$, and does the CPU/GPU forward+loss
throughput ratio $r$ determine where it applies?

\textbf{Hypothesis.} Near parity, dynamic self-scheduling~\cite{kruskal1985selfscheduling}
absorbs useful CPU work and reduces $T_{\text{grid}}$
against the matched GPU-only path. Far below parity,
GPU-only dominates or the result is unresolved. We also 
expect more headroom on the discrete-memory T4 than on the
unified-memory M4, since unified memory shares one bandwidth
budget between CPU and GPU
(Section~\ref{rq2-hybrid-scheduling-conditions}).

\textbf{Design.} Native $r$ on the available hardware does
not span the regimes where the ratio can change the outcome, so
the GPU slowdown factor sweeps $r$ to find the hybrid applicability threshold (Section~\ref{methods-scheduling}); the queue-based
scheduler sees only task-completion timing, so the sleep is
sufficient at this granularity. Each workload/platform pair
first records $r_{\text{native}}$; if $r_{\text{native}}<1$,
the slowdown factor $s=1/r_{\text{native}}$ targets parity,
otherwise $s=1$. At the achieved ratio, \texttt{hybrid} is
paired against \texttt{vanilla} with the same slowdown
applied to the GPU side. The GPU path remains
\texttt{vanilla}, so the measured effect belongs to the
scheduler axis.

Native rows and slowed rows answer different questions. The
native-$r$ row runs with no added latency, so it is the only
row claiming a speedup realizable on this hardware. The slowed
rows distort $r$ to locate where adding CPU workers stops
helping, without claiming a gain on
this machine. Table~\ref{tab:hybrid-ratio-check} keeps native
$r$ and achieved $r$ in separate columns so the two are not
conflated.

\subsection{Calibration Cache Amortization
(RQ3)}\label{experiment-3-cache}

\textbf{Question.} When the unit of evaluation changes from
one grid to a related-checkpoint session, can the
unconditional calibration step amortize via 
cached-policy reuse?

\textbf{Motivation.} A LossLens-style
session~\cite{xie2025losslens} evaluates several related
checkpoints in succession. Hybrid scheduling's applicability
varies by machine and workload (RQ2), so a session cannot commit to hybrid statically. It must
calibrate, and calibration is allowed to select GPU-only as
a valid outcome. The question is whether the one-time
calibration cost is recovered over the session. RQ3 is a
scoped measurement: it uses one favorable regime and a
single session measurement per condition, so it supports or
refutes amortization for that case and does not establish a
general amortization claim across other workloads or machines.

\textbf{Hypothesis.} Amortization holds where Experiment~2
gives positive hybrid evidence and the per-checkpoint
hybrid saving exceeds the calibration cost amortized over
$N=4$.

\textbf{Design.} RQ3 filters Experiment~2 for a
(workload, regime) with hybrid affinity, meaning a supported
hybrid win. It prefers a native regime, where the win is a
practical-hardware claim, and falls back to a slowed regime,
where the win is a controlled demonstration. The selected
regime's slowdown is inherited as a single operating point and
applied uniformly to calibration, \texttt{vanilla}, and
\texttt{cached\_composed}, so the comparison is fair within
that regime. Each session contains $N=4$ checkpoints from one
training trajectory, drawn at evenly spaced epochs (the exact
spacing recorded with the asset manifest). The GPU-side
configuration is \texttt{rq3\_config}, selected by rule: the
RQ1 candidate with the highest supported speedup at the
workload/platform pair, taking the smaller $K$ on ties for
\texttt{vmap}-bearing candidates, and falling back to baseline
when no RQ1 candidate has a supported improvement. Two
conditions are compared: \texttt{vanilla} run once per
checkpoint, and \texttt{cached\_composed}, which calibrates on
checkpoint $0$ and reuses the selected tuple for all
checkpoints under \texttt{rq3\_config}. Amortization here means
\texttt{cached\_composed} has session speedup $>1.0$ after
calibration time is included; the verdict is labeled native or
slowed, and break-even $N^{*}$ follows
Section~\ref{methods-calibration}.

Two upstream outcomes shape RQ3. If RQ1 gives no supported
improvement at the selected pair, \texttt{rq3\_config} falls
back to baseline (Section~\ref{methods-calibration}), so RQ3
still runs but over the GPU-only evaluation path. If
Experiment~2 shows hybrid affinity in no regime, native or
slowed, there is no regime to inherit, and RQ3 reports
amortization refuted without running a session because
calibration cannot amortize anywhere. Within a selected regime,
calibration may still resolve to GPU-only, which also refutes
amortization. RQ3 therefore results in either a regime where
reuse amortizes or evidence that it does not.
